% Part: second-order-logic
% Chapter: syntax-and-semantics
% Section: size-of-domain

\documentclass[../../../include/open-logic-section]{subfiles}

\begin{document}

\olfileid{sol}{syn}{siz}

\olsection{అనంతమైన, \usetoken{S}{enumerable}
  \usetoken{P}{domain}లను వివరించడం}

సమితి~$M$ (డెడెకిండ్) అనంతం అవుతుంది అప్పుడే, అప్పుడు మాత్రమే
\tetoken{సంగ్రస్త}{!!{surjective}} కాని \tetoken{ఒకటి-ఒకటి}{!!a{injective}} ప్రమేయం~$f\colon M \to M$ ఒకటి
ఉంటుంది; అంటే $\ran{f} \neq M$. మొదటిస్థాయి తర్కంలో, ఒక ఏకస్థాన
\tetoken{ప్రమేయ సంకేతం}{!!{function}}~$f$ను పరిగణించి, \tetoken{ఒక నిర్మాణం}{!!a{structure}}~$\Struct{M}$లో దానికి
నిర్దేశించిన ప్రమేయం~$\Assign{f}{M}$ \tetoken{ఒకటి-ఒకటి}{!!{injective}} అనీ,
$\ran{f} \neq \Domain{M}$ అనీ ఇలా చెప్పవచ్చు:
\[
\lforall[x][\lforall[y][(\eq[f(x)][f(y)] \to \eq[x][y])]] \land
\lexists[y][\lforall[x][\eq/[y][f(x)]]].
\]
$\Struct{M}$ ఈ \tetoken{వాక్యాన్ని}{!!{sentence}} సంతృప్తిపరిస్తే, $\Assign{f}{M}:
\Domain{M} \to \Domain{M}$ \tetoken{ఒకటి-ఒకటి}{!!{injective}}; కాబట్టి $\Domain{M}$
అనంతం కావాలి. $\Domain{M}$ అనంతమైతే, అందువల్ల అటువంటి ప్రమేయం
ఉంటే, $\Assign{f}{M}$ను ఆ ప్రమేయంగా తీసుకోవచ్చు; అప్పుడు $\Struct{M}$
ఆ \tetoken{వాక్యాన్ని}{!!{sentence}} సంతృప్తిపరుస్తుంది. అయితే, ఈ ప్రయోజనం కోసం ఉపయోగించే
తార్కికేతర సంకేతం~$f$ మన భాషలో ఉండాలని ఇది కోరుతుంది. ద్వితీయ-స్థాయి
తర్కంలో, అటువంటి ప్రమేయం \emph{ఉందని} నేరుగా చెప్పవచ్చు. అప్పుడు~$f$
అవసరం ఉండదు; శుద్ధ ద్వితీయ-స్థాయి తర్కంలో కింది \tetoken{వాక్యం}{!!{sentence}} లభిస్తుంది:
\[
\fn{Inf} \ident \lexists[u][(\lforall[x][\lforall[y][(\eq[u(x)][u(y)]
      \lif \eq[x][y])]] \land \lexists[y][\lforall[x][\eq/[y][u(x)]]])].
\]
$\Sat{M}{\fn{Inf}}$ అవుతుంది అప్పుడే, అప్పుడు మాత్రమే $\Domain{M}$
అనంతం. తరువాత $\fn{Fin} \ident \lnot \fn{Inf}$ అని నిర్వచించవచ్చు;
$\Sat{M}{\fn{Fin}}$ అవుతుంది అప్పుడే, అప్పుడు మాత్రమే $\Domain{M}$
పరిమితం. శుద్ధ మొదటిస్థాయి తర్కంలోని ఏ ఒక్క \tetoken{వాక్యమూ}{!!{sentence}}
\tetoken{వ్యక్తి క్షేత్రం}{!!{domain}} అనంతమని వ్యక్తం చేయలేదు; అయితే వాటి అనంత సమితి ఒకటి
వ్యక్తం చేయగలదు. ఒక \tetoken{నిర్మాణంలో}{!!a{structure}} దాని వ్యక్తి క్షేత్రం పరిమితమైనప్పుడు,
అప్పుడు మాత్రమే సంతృప్తమయ్యే శుద్ధ మొదటిస్థాయి తర్క
\tetoken{వాక్యాల}{!!{sentence}s} సమితి ఏదీ లేదు.

\begin{prop}
$\Sat{M}{\fn{Inf}}$ అవుతుంది అప్పుడే, అప్పుడు మాత్రమే $\Domain{M}$
అనంతం.
\end{prop}

\begin{proof}
$\Sat{M}{\fn{Inf}}$ అవుతుంది అప్పుడే, అప్పుడు మాత్రమే ఏదో ఒక~$s$కు
  $\Sat{M}{\lforall[x][\lforall[y][(\eq[u(x)][u(y)] \lif \eq[x][y])]]
    \land \lexists[y][\lforall[x][\eq/[y][u(x)]]]}[s]$.
  అది జరిగితే, $s(u)$ \tetoken{ఒకటి-ఒకటి}{!!a{injective}} ప్రమేయం; ఏదో ఒక $y
  \in \Domain{M}$, ~$s(u)$ వ్యాప్తిలో ఉండదు. విపర్యయంగా,
  $\ran{f} \neq \Domain{M}$ అయ్యే \tetoken{ఒకటి-ఒకటి}{!!a{injective}} ప్రమేయం
  $f\colon \Domain{M} \to \Domain{M}$ ఉంటే, $s(u) = f$ అటువంటి చర
  నిర్దేశం.
\end{proof}

సమితి~$M$ యొక్క \tetoken{మూలకాల}{!!{element}s} జాబితా
\[
m_0, m_1, m_2, \dots
\]
పునరుక్తులు లేకుండా (అయితే బహుశా పరిమితంగా) ఉంటే, ఆ సమితిని
\tetoken{లెక్కించదగినది}{!!{enumerable}} అంటాం. \tetoken{ఒక మూలకం}{!!a{element}}~$z \in M$, ఒక ప్రమేయం
$f\colon M \to M$ ఉండి, $z$, $f(z)$, $f(f(z))$, \dots, అనేవి
$M$ యొక్క అన్ని \tetoken{మూలకాలు}{!!{element}s} అయినప్పుడు, అప్పుడు మాత్రమే అటువంటి
జాబితా ఉంటుంది. జాబితా ఉంటే, $z = m_0$, $f(m_k) = m_{k+1}$గా
(లేదా జాబితాలో $m_k$ చివరి \tetoken{మూలకం}{!!{element}} అయితే $f(m_k) = m_k$గా)
తీసుకుంటే, కావలసిన \tetoken{మూలకం}{!!{element}}, ప్రమేయం లభిస్తాయి. మరోవైపు, అటువంటి
$z$,~$f$ ఉంటే, $z$, $f(z)$, $f(f(z))$, \dots, అనేది~$M$ యొక్క
జాబితా; కాబట్టి $M$ \tetoken{లెక్కించదగినది}{!!{enumerable}}. ఒక \tetoken{నిర్మాణంలో}{!!a{structure}}, ఆ
\tetoken{నిర్మాణం}{!!{structure}} \tetoken{లెక్కించదగినదైనప్పుడు}{!!{enumerable}}, అప్పుడు మాత్రమే సత్యమయ్యే \tetoken{వాక్యాన్ని}{!!{sentence}}
ఉత్పత్తి చేయడానికి ద్వితీయ-స్థాయి తర్కంలో $z$,~$f$ల అస్తిత్వాన్ని
ఇలా వ్యక్తం చేయవచ్చు:
\[
\fn{Count} \ident
\lexists[z][\lexists[u][\lforall[X][((X(z) \land
      \lforall[x][(X(x) \lif X(u(x)))]) \lif \lforall[x][X(x)])]]]
\]

\begin{prop}
$\Sat{M}{\fn{Count}}$ అవుతుంది అప్పుడే, అప్పుడు మాత్రమే $\Domain{M}$
\tetoken{లెక్కించదగినది}{!!{enumerable}}.
\end{prop}

\begin{proof}
$\Domain{M}$ \tetoken{లెక్కించదగినది}{!!{enumerable}} అనుకుందాం; $m_0$, $m_1$, \dots,
ఒక జాబితా అనుకుందాం. పునరుక్తులను తొలగించడం ద్వారా ఏ $m_k$ కూడా
రెండుసార్లు రాకుండా హామీ ఇవ్వవచ్చు. $f(m_k) = m_{k+1}$గా
నిర్వచించి, $s(z) = m_0$, $s(u) = f$గా తీసుకుందాం. కిందిది
సత్యమని చూపుతాం:
\[
\Sat{M}{\lforall[X][((X(z) \land \lforall[x][(X(x) \lif X(u(x)))])
    \lif \lforall[x][X(x)])]}[s]
\]
$S \subseteq \Domain{M}$ యాదృచ్ఛికం అనుకుందాం. ఇంకా,
$\Sat{M}{(X(z) \land \lforall[x][(X(x) \lif
X(u(x)))])}[\Subst{s}{S}{X}]$ అనుకుందాం. అప్పుడు
$\Subst{s}{S}{X}(z) \in S$; $x \in S$ అయినప్పుడల్లా
$(\Subst{s}{S}{X}(u))(x) \in S$ కూడా. మరోలా చెప్పాలంటే,
$\varAssign{\Subst{s}{S}{X}}{s}{X}$ కాబట్టి $m_0 \in S$; $x \in S$
అయితే $f(x) \in S$. అందువల్ల $m_0 \in S$, $m_1 = f(m_0) \in S$,
$m_2 = f(f(m_0)) \in S$, ఇలాగే మిగతావీ. కాబట్టి
$S = \Domain{M}$; అందువల్ల
$\Sat{M}{\lforall[x][X(x)]}[\Subst{s}{S}{X}]$. $S \subseteq
\Domain{M}$ యాదృచ్ఛికం కాబట్టి నిరూపణ పూర్తయింది:
$\Sat{M}{\fn{Count}}$.
\sourcecorrection{OLTESOLSYN-005}{ముందుకు దిశలో ఆంగ్ల మూలం నిర్మాణం $M$, దాని యాదృచ్ఛిక ఉపసమితి రెండింటికీ $M$నే వాడింది. నిర్మాణాన్ని $M$గానే ఉంచి ఉపసమితిని $S$గా మార్చి, ఆ పేరును సంబంధ-చర ప్రతిస్థాపనలన్నింటిలో అనుసరించాం.}

ఇప్పుడు $\Sat{M}{\fn{Count}}$, అంటే ఏదో ఒక~$s$కు
\[
\Sat{M}{\lforall[X][((X(z) \land \lforall[x][(X(x) \lif X(u(x)))])
    \lif \lforall[x][X(x)])]}[s]
\]
అని అనుకుందాం. $m = s(z)$, $f = s(u)$గా తీసుకొని,
$S = \{m, f(m), f(f(m)), \dots\}$ను పరిగణించండి. ఇలా నిర్వచించిన
$S$ స్పష్టంగా \tetoken{లెక్కించదగినది}{!!{enumerable}}. అప్పుడు ఊహ ప్రకారం
\[
\Sat{M}{(X(z) \land \lforall[x][(X(x) \lif X(u(x)))])
    \lif \lforall[x][X(x)]}[\Subst{s}{S}{X}]
\]
అంతేకాక, $S \ni m = \Subst{s}{S}{X}(z)$ కాబట్టి
$\Sat{M}{X(z)}[\Subst{s}{S}{X}]$; $x \in S$ అయినప్పుడల్లా
$f(x) \in S$ కూడా కాబట్టి
$\Sat{M}{\lforall[x][(X(x) \lif
X(u(x))) ]}[\Subst{s}{S}{X}]$. పూర్వపక్షం, సోపాధికం రెండూ
సంతృప్తమయ్యాయి కాబట్టి, పర్యవసానం కూడా సంతృప్తమవాలి:
$\Sat{M}{\lforall[x][X(x)]}[\Subst{s}{S}{X}]$. కానీ దాని అర్థం
$S = \Domain{M}$. కాబట్టి నిర్వచనం ప్రకారం $S$ లెక్కించదగినది కనుక
$\Domain{M}$ కూడా \tetoken{లెక్కించదగినది}{!!{enumerable}}.
\sourcecorrection{OLTESOLSYN-006}{వెనుకకు దిశలో ఆంగ్ల మూలం నిర్మాణం $M$, దాని ఒక నిర్దిష్ట ఉపసమితి రెండింటికీ $M$నే వాడింది. నిర్మాణాన్ని $M$గానే ఉంచి, $m$ యొక్క $f$-పునరావృతాల సమితిని $S$గా మార్చి, ఆ పేరును ప్రతిస్థాపనలన్నింటిలో అనుసరించాం.}
\end{proof}

\begin{prob}
\tetoken{వాక్యం}{!!{sentence}} $\fn{Inf} \land \fn{Count}$ సరిగ్గా
\tetoken{అనంతంగా లెక్కించదగిన}{!!{denumerable}} వ్యక్తి క్షేత్రాలన్నింటిలో, వాటిలో మాత్రమే
సత్యం. $\fn{Count}$ నిర్వచనాన్ని సవరించి, వ్యక్తి క్షేత్రం
\tetoken{అనంతంగా లెక్కించదగినది}{!!{denumerable}} అని నేరుగా వ్యక్తం చేసే వేరొక
\tetoken{వాక్యంగా}{!!{sentence}} దాన్ని మార్చండి; అది అలా చేస్తుందని నిరూపించండి.
\end{prob}

\end{document}
